% ===================================================================
\chapter{Hardware Implementation}
\label{ch:hardware}

Both algorithms were implemented as a compressor core and a matching
decompressor core in C++ for Vitis HLS, and built into a test system on an AMD
KCU116 evaluation board. This chapter describes the cores, the measures that
keep them exactly lossless on any input, and the board system that runs and
measures them.

\section{Common interface and bit packing}
\label{sec:interface}

All four cores share one interface. A compressor accepts AXI4-Stream input at
\num{128} bits per beat --- four receive channels, each a 16-bit real and a
16-bit imaginary part --- and emits a compressed AXI4-Stream at \num{256} bits
per beat, with TLAST marking the end of the frame. A decompressor does the
reverse. Frame dimensions (range bins, ramps, receive channels and $\nTx$)
arrive on an AXI4-Lite control interface, so there is no header in the
bitstream.

Both compressors share the bit packer. Codewords accumulate in a 512-bit
register, least-significant-bit first, and a 256-bit word is written out
whenever at least 256 bits are waiting, with one final partial word at the end
of the frame carrying TLAST. If the frame ends exactly on a word boundary there
is no partial word, so an all-zero word is sent to carry TLAST; without it the
DMA and the decoder would wait for an end of frame that never comes. No
ColoRadar frame ends on a boundary --- the case was found by the synthetic
edge-case frames of Section~\ref{sec:dechw}, and the fix leaves the output for
all 50 ColoRadar frames byte-identical. Flushing per frame rather than per ramp avoids wasting up to 255
bits per ramp, about \SI{2.9}{\percent} of the ratio. The Huffman codewords are
stored bit-reversed relative to their usual most-significant-bit-first form,
because the packer is least-significant-bit first; the reversed code set is
still prefix-free, so a decoder can match candidates from the low end of its
buffer with no ambiguity. Compressor and decompressor read the same
sixteen-entry table. Each core was built with both dictionaries of
Section~\ref{sec:dict}; the two builds differ only in the contents of this table,
and have the same schedule and timing. Both were run on the board
(Chapter~\ref{ch:results}); the compressors' cycle counts agree to within 13
cycles per frame.

\section{DRHE compressor}
\label{sec:drheimpl}

\subsection{State and parallelism}

DRHE holds five \texttt{float} state arrays --- the magnitude prediction, the
previous magnitude, the phase prediction and the previous two phases --- each
indexed by (transmitter, receive channel, range bin). The transmitter index is
the ramp number modulo $\nTx$, which is all it takes to put the predictor on the
per-transmitter axis of Section~\ref{sec:predaxis}: indexing by transmitter
changes \emph{which} address is read, not \emph{when}, so the schedule is the
same as for a lag-1 design. Each transmitter's history starts on its own first
ramp, so the state is reset while $m < \nTx$. Frames are independent:
nothing carries over from one frame to the next, so a frame lost in transit
cannot corrupt its successors.

The channel dimension is completely partitioned, so the four receive channels
are four parallel copies of the datapath inside one pipelined iteration of the
sample loop. The state occupies 40 block RAMs (36\,Kb).

\subsection{Where the logic goes}

Per sample and per channel, DRHE evaluates \texttt{sqrt} and \texttt{atan2} to
convert the reconstructed sample to polar form, and \texttt{sin} and
\texttt{cos} to convert the prediction back --- four floating-point
transcendental functions, for each of four channels, on every one of the
\num{24576} beats of a frame. These operators dominate the design: they account
for about two thirds of its LUTs and all 132 of its DSP slices, and the two
arctangent units alone are about a third of the logic. Kiem's design uses fixed
point for the same functions; Chapter~\ref{ch:hwinterpret} measures what that choice costs.

\subsection{Initiation interval}

DRHE runs at $\mathrm{II} = 2$, and the reason has nothing to do with
arithmetic. Each state array is partitioned by channel, so each channel's
history is one dual-port block RAM. Every iteration reads a value and writes
the updated one --- two accesses, which two ports can serve. On the first ramps
of a frame the iteration also writes a reset value: three accesses against two
ports, which forces the tool to serialise, and the loop settles at
$\mathrm{II} = 2$ for the whole frame. Reading the state as zero on the first
ramps instead of writing zeros to it reaches $\mathrm{II} = 1$. It was not
adopted, because it doubles the floating-point arithmetic (units can no longer
be shared between two channels) and the sensor does not need the throughput:
the $\mathrm{II} = 2$ design already runs many times faster than the radar
produces frames (Section~\ref{sec:boardtiming}).

\section{LPC compressor}
\label{sec:lpcimpl}

\subsection{Three passes}

LPC cannot stream, so its structure is explicitly multi-pass
(Figure~\ref{fig:lpcblock}):

\begin{enumerate}
\item \textbf{Ingest.} Read the whole frame into a buffer.
\item \textbf{Correlate, solve and residual}, once per (channel, bin): compute
$R_0$, $R_1$ and $R_2$ at lags $0$, $\nTx$ and $2\nTx$ from the buffer, solve
Equation~\eqref{eq:yw}, then walk the ramps again computing residuals and
writing them back into the buffer \emph{in place}.
\item \textbf{Emit.} Write the coefficient table, then the coded residuals.
\end{enumerate}

\noindent
Rewriting the buffer in place means one \SI{3.0}{\mebi bit} memory serves both
purposes rather than two.

\begin{figure}[t]
\centering
\begin{tikzpicture}[font=\small,
  b/.style={draw=slate,fill=white,rounded corners=2pt,minimum width=30mm,minimum height=9mm,font=\scriptsize,align=center},
  h/.style={b,draw=accent,line width=0.9pt,fill=lightbg},
  m/.style={draw=slate,fill=warm!12,rounded corners=2pt,minimum width=18mm,minimum height=22mm,font=\scriptsize,align=center},
  a/.style={-{Latex[length=1.5mm]},slate}]
\node[b]  (in)  at (0,0)    {AXI in\\128\,b/beat};
\node[h]  (p1)  at (3.3,0)  {\textbf{pass 1}\\ingest};
\node[m]  (buf) at (6.6,0)  {frame\\buffer\\\SI{3.0}{\mebi bit}};
\node[h]  (p2)  at (9.9,0)  {\textbf{pass 2}, per bin\\correlate, solve,\\residual (in place)};
\node[h,minimum width=44mm]  (p3)  at (6.6,-2.5) {\textbf{pass 3} --- emit coefficients,\\then residuals in raster order};
\node[b]  (out) at (0,-2.4) {AXI out\\256\,b/beat};
\draw[a] (in) -- (p1);
\draw[a] (p1) -- (buf);
\draw[a] (buf) -- (p2);
\draw[a] (p2.south) -- (9.9,-2.5) -- (p3.east);
\draw[a] (p3) -- (out);
\node[font=\scriptsize,text=slate,align=center] at (9.9,1.15) {$R_0,R_1,R_2$ at lags $0,\nTx,2\nTx$};
\end{tikzpicture}
\caption[LPC three-pass architecture]{The LPC compressor. Because the
Yule--Walker coefficients depend on the whole frame, the compressor must buffer
it; each (channel, bin) is then correlated, solved and turned into residuals,
which overwrite the buffer in place so that only one memory is needed. Emitting residuals in the original raster order is what lets
the \emph{decompressor} avoid a frame buffer entirely.}
\label{fig:lpcblock}
\end{figure}

\subsection{Reaching back $\nTx$ and $2\nTx$ ramps}

Both the correlation and the residual loops need $x[m - \nTx]$ and
$x[m - 2\nTx]$. A shift register of depth $2\nTx = 24$ would move 24 entries
every iteration. Instead a \emph{ring} of exactly $2\nTx$ entries gives both
taps for free. At ramp $m$, let $k = m \bmod 2\nTx$. Slot $k$ still holds the
value written at ramp $m - 2\nTx$, because it is about to be overwritten and has
not been touched since; and slot $(k + \nTx) \bmod 2\nTx$ holds the value from
ramp $m - \nTx$. So:
\begin{center}
\begin{tabular}{@{}ll@{}}
read $x[m-2\nTx]$ & from slot $k$ \\
read $x[m-\nTx]$  & from slot $(k+\nTx) \bmod 2\nTx$ \\
write $x[m]$      & to slot $k$ \\
\end{tabular}
\end{center}
Nothing is shifted; the cost is $2\nTx$ registers and two multiplexers. Every
loop in the design reaches $\mathrm{II} = 1$.

\subsection{Where the cycles go}

A frame takes \num{325366} cycles on the board. Ingest and emit are
\num{24576} cycles each, one per beat. The remaining \num{276000} or so are the
per-bin pass, which runs two short pipelined loops --- correlation and residual
--- for each of the 512 (channel, bin) pairs. A pipelined loop pays its
pipeline depth once on entry before reaching steady state, about \num{71} cycles
for the correlation loop and \num{56} for the residual loop, so 512 invocations
add roughly $512 \times (71 + 56) \approx \num{65000}$ cycles of fill that a
trip-count model does not show. Merging the two loops, or processing several
bins per invocation, would recover most of it; overlapping the three passes
across frames would bring the time per frame down further. Neither was
necessary for this thesis, and both are listed as future work.

\subsection{Fixed-point choices}

Coefficients are quantised to \texttt{ap\_fixed<16,2>} for $a_1$ (range
$[-2,2)$) and \texttt{ap\_fixed<16,1>} for $a_2$ (range $[-1,1)$), which follows
from the stability conditions of Equation~\eqref{eq:stability}. Over 50 frames
this quantisation costs nothing measurable: the residuals are identical to the
bit to those of double-precision coefficients.

The autocorrelations are accumulated as exact 64-bit integers rather than in
floating point. $R_0^2 - R_1^2$ is a difference of two nearly equal quantities
whenever the data is strongly correlated --- precisely the regime the predictor
exists for --- so computing it in floating point would lose most of its
significant digits exactly when accuracy matters most. The numerators and the
determinant are formed exactly as 128-bit integers before the single division.

\subsection{An asymmetry that favours the decoder}

Pass 3 emits the coefficient table first --- the decoder needs all of it before
it can reconstruct anything --- and then the residuals \emph{in the raster
order the samples arrived in}, not in the (channel, bin, ramp) order pass 2
produced them. The decompressor therefore needs no frame buffer at all: it holds
the coefficient table and the last $2\nTx$ reconstructed samples of each
(channel, bin), and writes each output sample as its residual arrives. The
reordering is paid for once, by the side that already had to buffer. LPC's
non-causality is a compressor-side cost only.

\section{Exactly lossless on any input}
\label{sec:robust}

Real radar data exercises only the lower part of the entropy coder's range
(Section~\ref{sec:edgecases}). Two situations that real data never produces, but synthetic input does
(Section~\ref{sec:edgecases}), are handled explicitly so that both designs are
lossless on \emph{any} input.

\subsection{The one value the code cannot represent}

Category 15 of the $S_4$ code covers $|v| \in [16384, 32767]$: counting both
signs, exactly the $2^{15}$ patterns its fifteen payload bits can hold. But
16-bit arithmetic can produce $v = -32768$, for which
Equation~\eqref{eq:append} gives $A(-32768) = -1$, which truncated to fifteen
bits is \texttt{0x7FFF} --- the same payload as $v = +32767$. The two would be
indistinguishable on the wire.

Three remedies were considered: extend the code with a category 16, which
changes the bitstream format; reject the input; or saturate $-32768$ to
$-32767$. Saturation was chosen. It costs at most one least-significant bit on
the affected sample, leaves the format untouched, and never fires on real data.

Saturation alone is not enough, though, and this is where the closed-loop
prediction of Section~\ref{sec:closedloop} becomes concrete. If the encoder
drove its predictor from the \emph{original} sample while the decoder drives
its own from the \emph{reconstructed} one, a single saturated sample would make
them diverge, and because DRHE's filters are recursive that divergence would
carry through every remaining ramp. The encoder therefore reconstructs exactly
what the decoder will see and updates its state from that:

\begin{verbatim}
int diff_re = wrap_int16(re - pred_re);
if (diff_re == -32768) diff_re = -32767;      // saturate
int recon_re = wrap_int16(diff_re + pred_re); // what the decoder will see
float curr_mag   = hls::sqrt(recon_re*recon_re + recon_im*recon_im);
float curr_phase = hls::atan2(recon_im, recon_re);   // state from recon
\end{verbatim}

\noindent
On real data the reconstruction always equals the original, so the closed loop
changes no output bit. Its cost is in the schedule: the reconstruction sits
ahead of the polar conversion on the dependency chain, and accounts for 26 of
the pipeline's 58 stages.

\subsection{Wide-integer conversion in the coefficient solver}

The LPC solver forms its numerators and determinant as 128-bit integers. The
conversion of such a value directly to \texttt{float} is not correct above
$2^{64}$ in this toolchain. All three operands are therefore normalised by one
common arithmetic right shift into 62 bits before the division, with the shift
chosen from the largest of them:
\begin{equation}
k = \max\bigl(\mathrm{msb}|D|,\ \mathrm{msb}|N_1|,\ \mathrm{msb}|N_2|\bigr) - 61,
\end{equation}
clamped at zero. Because numerator and denominator are shifted alike, the
quotient is unchanged. On ColoRadar the correlations reach only 44 bits, so
the shift is zero; on full-scale synthetic input it is what keeps the
coefficients correct.

\section{Decompressors}
\label{sec:dechw}

A decompressor exists for each algorithm, for two reasons. It is what makes
bit-exact verification possible --- every frame is compressed, decompressed and
compared --- and it completes the pipeline a vehicle would need, with the
encoder at the sensor and the decoder at the central processor. Each
decompressor is built twice from the same C++ source: as an FPGA core, and as a
PC program compiled through Vitis C-simulation, which uses the same
bit-accurate models of the floating-point functions.

\subsection{Why decoding is the harder half}

Entropy decoding is inherently serial: the length of a codeword is known only
once it has been matched, and the next codeword starts where it ends. A DRHE
decoder that keeps pace with the compressor --- eight values (four channels,
real and imaginary) per two clocks --- needs eight chained codeword matches and
very wide variable shifters per cycle, plus four copies of the floating-point
predictor. HLS estimates such a design at about \num{182000} LUTs. HLS
over-estimates (by about three times for the serial decoder below), but even a
third of that would be four times the serial decoder's logic. Keeping pace is
also unnecessary: a decoder that takes \SI{2}{\milli\second} per frame is a
hundred times faster than the radar produces frames. The DRHE decoder was
therefore designed to be small rather than fast.

\subsection{DRHE: a two-stage serial decoder}

The DRHE decompressor is therefore a dataflow pair of stages connected by a
stream:

\begin{itemize}
\item \textbf{Unpack} decodes one value per clock. It cuts a 29-bit window from
two 256-bit input words, matches all sixteen codewords against it in parallel,
extracts the payload, and reads the next input word just in time, stopping at
TLAST.
\item \textbf{Rebuild} runs one floating-point predictor that handles one
receive channel per step at $\mathrm{II} = 2$ (limited by the ports of the
sine/cosine tables), keeps its state in block RAM, and uses exactly the same
expressions as the compressor so that its arithmetic matches bit for bit.
\end{itemize}

\noindent
The two stages are balanced at about \num{196608} clocks per frame, roughly
\SI{2}{\milli\second} at \SI{100}{\mega\hertz}; on the board it takes
\num{196881} (Section~\ref{sec:boardtiming}). HLS estimated \num{42846} LUTs; after
placement and routing the core uses \num{14155}.

Bit-for-bit agreement of the floating-point arithmetic is not optional for
DRHE. The predictor feeds back on itself, so a decoder whose \texttt{atan2} or
\texttt{sin} differs from the encoder's in the last bit will drift. A
double-precision Python decoder written as a cross-check is exact for LPC but
drifts by 1 to 6 least-significant bits on 8 of the 50 DRHE frames; the
decoder built from the HLS source, which the FPGA core shares, is exact on all
of them.

\subsection{LPC: the same two stages}

The LPC decompressor has the same structure as the DRHE decoder, a dataflow
pair of stages joined by a stream:

\begin{itemize}
\item \textbf{Unpack} first reads the coefficient table, one 64-bit entry per
clock, and passes it on. It then decodes one residual per clock with the DRHE
decoder's matcher: a 29-bit window cut from two 256-bit input words, all sixteen
codewords compared with it at once, the payload extracted, and the next input
word read just in time, stopping at TLAST.
\item \textbf{Rebuild} stores the coefficients and then reconstructs one complex
sample per clock: two fixed-point multiply--accumulates for the real part and
two for the imaginary part, the same rounding, and the same 16-bit wrap as the
compressor. It holds the last $2\nTx$ reconstructed values of each (channel,
bin) in a ring, as the compressor does. The rings are in distributed (LUT) RAM,
split so that both past values are read and the new one written in the same
clock.
\end{itemize}

\noindent
Unpack is the slower stage, at one real value per clock, so a frame takes about
\num{196608} clocks plus 512 for the coefficient table --- about
\SI{2}{\milli\second}, the same as DRHE. Every loop is scheduled at
$\mathrm{II} = 1$, and the decoder needs no frame buffer and no block RAM.
After synthesis it uses \num{15359} LUTs (just over half of them as distributed
RAM for the rings), \num{7212} flip-flops and 23 DSP slices.

The first LPC decoder built was simpler: it tried the sixteen codewords one
after another in category order and stopped at the first match. It was exact,
but its speed depended on the data. On the board it took \SI{15.4}{\milli\second}
per frame, 7.6 to 8.3 clocks per value, and \num{20034} LUTs after routing ---
slower and larger than the parallel matcher, which therefore replaced it. In
C-simulation the new decoder's output is bit-identical to the first one's on
all 50 ColoRadar frames and on eight synthetic edge-case frames: all zeros,
constant $+32767$ and $-32768$, alternating full scale, uniform random values,
small noise, a random walk along the ramps, and small noise with full-scale
spikes.

\section{The board test system}
\label{sec:system}

The cores run in a complete system on an AMD KCU116 evaluation board, whose
Kintex UltraScale+ \texttt{XCKU5P} provides \num{216960} LUTs, \num{433920}
flip-flops, 480 block RAMs (36\,Kb) and \num{1824} DSP slices
(Figure~\ref{fig:system}).

\begin{figure}[t]
\centering
\begin{tikzpicture}[font=\small,x=0.88cm,
  b/.style={draw=slate,fill=white,rounded corners=2pt,minimum height=8mm,inner xsep=4pt,font=\scriptsize,align=center},
  h/.style={b,draw=accent,line width=0.9pt,fill=lightbg},
  a/.style={{Latex[length=1.4mm]}-{Latex[length=1.4mm]},slate}]
\node[b,fill=lightbg,minimum width=20mm] (pc) at (-3.4,1.7) {host PC\\(JTAG, UART)};
\node[b,minimum width=24mm] (mb) at (0,1.7) {MicroBlaze\\128\,KB local RAM};
\node[b,minimum width=24mm] (io) at (0,0.1) {UART, timer,\\LEDs};
\node[b,minimum width=20mm,minimum height=32mm] (ic) at (3.4,0.9) {AXI\\SmartConnect};
\node[b,minimum width=24mm] (ddr) at (7.0,2.5) {DDR4-2400, 1\,GB\\frames and results};
\node[b,minimum width=18mm] (dc) at (7.0,0.9) {DMA \texttt{dma\_c}};
\node[b,minimum width=18mm] (dd) at (7.0,-0.7) {DMA \texttt{dma\_d}};
\node[h,minimum width=28mm] (cc) at (10.6,0.9) {compressor\\128\,b in, 256\,b out};
\node[h,minimum width=28mm] (dk) at (10.6,-0.7) {decompressor\\256\,b in, 128\,b out};
\draw[a] (pc) -- (mb);
\draw[a] (mb) -- (mb -| ic.west);
\draw[a] (io) -- (io -| ic.west);
\draw[a] (ic.east |- ddr) -- (ddr);
\draw[a] (ic.east |- dc) -- (dc);
\draw[a] (ic.east |- dd) -- (dd);
\draw[a] (dc) -- (cc);
\draw[a] (dd) -- (dk);
\node[font=\scriptsize,text=slate] at (10.6,-1.45) {\texttt{\_full} builds only};
\end{tikzpicture}
\caption[The KCU116 test system]{The board test system. The host places the
frames in DDR4 over JTAG; the MicroBlaze starts the DMAs, times each frame and
writes a result record; the host then reads the results and compressed frames
back. Everything runs at
\SI{100}{\mega\hertz}.}
\label{fig:system}
\end{figure}

\paragraph{Hardware.} A MicroBlaze processor with \SI{128}{\kilo\byte} of local
memory controls the system. A DDR4-2400 controller gives access to \SI{1}{\giga\byte}
of board memory, which holds the input frames, the compressed frames, the
reconstructed frames and the results. An AXI DMA (\texttt{dma\_c}) streams
each frame from DDR4 through the compressor and writes the compressed stream
back. In the \texttt{\_full} build a second DMA (\texttt{dma\_d}) streams the
compressed frame through the decompressor and writes the reconstruction back,
so the board compresses, decompresses and checks every frame by itself. A UART,
a hardware timer and LEDs complete the system, all on one AXI interconnect at
\SI{100}{\mega\hertz}. Each algorithm therefore has two bitstreams: a
compress-only build and a \texttt{\_full} build with the decompressor.

\paragraph{Firmware.} For each frame the MicroBlaze configures the core's frame
dimensions, starts the DMA transfers, and reads a hardware timer at the start,
when the DMA's read channel reports the last input word read, and when its
write channel reports the last output word written. From these it
records the cycles per frame and the end-to-end latency of
Section~\ref{sec:metrics}, together with the compressed length; the
\texttt{\_full} firmware also records the decompressor's cycles and latency and
compares the reconstruction with the original word by word. Each frame's
results are a 16-word record in DDR4.

\paragraph{Host.} A script on the PC (\texttt{run\_board.bat}, driving
\texttt{xsdb}) programs the FPGA, loads the frames into DDR4 over JTAG, runs the
firmware, and reads back the result records and compressed frames. For the
compress-only build the PC then decompresses every frame with the decompressor
built from the HLS source. Finally \texttt{compare\_results.py} checks every
compressed frame byte for byte against the C-simulation output, checks every
reconstruction against the original frame, and tabulates compression ratio,
cycles and latency.

\section{Building and timing closure}
\label{sec:timing}

Each core is exported from Vitis HLS as an IP block and placed in a Vivado block
design with the system above; a script builds each of the four bitstreams. HLS
estimated both compressors below the \SI{100}{\mega\hertz} target. Its estimate
is the tool's own model of operator delays against a budget that reserves
about a quarter of the period for clock uncertainty, not a timing analysis of a
real circuit, and once placed and routed all four systems meet
\SI{100}{\mega\hertz} with positive slack (Table~\ref{tab:builds}). Two system
details were needed for that: a false path from the DDR4 calibration-complete
signal to the processor reset, and the DDR4 interface at \SI{2400}{MT/s}.

\begin{table}[t]
\centering
\caption[The four board builds]{The four board builds, whole system (processor,
DDR4 controller, DMAs and cores), post-route on the \texttt{XCKU5P} at
\SI{100}{\mega\hertz}, with the re-derived dictionaries; the original-dictionary
builds also meet timing. BRAM is in 36\,Kb blocks. \pending{LPC \texttt{\_full}
row: estimate for the parallel decoder; replace with the post-route report}}
\label{tab:builds}
\small
\begin{tabular}{@{}lrrrr@{}}
\toprule
Build & LUT & BRAM & DSP & WNS (ns)\\
\midrule
DRHE, compress only    & \num{61011} (\SI{28}{\percent}) & 104.5 & 138 & $+0.243$\\
LPC, compress only     & \num{59203} (\SI{27}{\percent}) & 162.5 & 164 & $+0.174$\\
DRHE, \texttt{\_full}  & \num{81573} (\SI{38}{\percent}) & 144   & 184 & $+0.341$\\
LPC, \texttt{\_full}   & \tbc{$\approx$\num{81000}} (\tbc{\SI{37}{\percent}}) & \tbc{169.5} & \tbc{187} & \tbc{$>0$}\\
\bottomrule
\end{tabular}
\end{table}
